\input{../tex-inputs/latex-leseansicht-vorspann}

\section[Gustav Schwarzkopf an Arthur Schnitzler, 6. 12. 1896]{L04286 Gustav Schwarzkopf an Arthur Schnitzler, 6. 12. 1896}
\nopagebreak\mylabel{L04286v}
\rehead{ }\normalsize\beginnumbering\briefempfaengerindex{Schnitzler, Arthur@\textsc{Schnitzler, Arthur}!zzzSchwarzkopf, Gustav@\emph{von Gustav Schwarzkopf}!1896-12-061@{6. 12. 1896}|(be}
\toendnotes[C]{\smallbreak\pagebreak[2]}
\correspDesc{Versand  durch Gustav Schwarzkopf am 6. 12. 1896 in Wien
\newline{}Erhalt  durch Arthur Schnitzler im Zeitraum [6. 12. 1896 – 9. 12. 1896?] in Wien}\toendnotes[C]{\smallbreak}
\Standort{CUL, Schnitzler, B 96a.}
\physDesc{Briefkarte, 400 Zeichen
\newline{}Handschrift: schwarze Tinte, deutsche Kurrent}\toendnotes[C]{\smallbreak}
\pstart
           \raggedleft{}{\pb}6. 12. 96\pend
           \vspace{0.5em}
\pstart
           Lieber Arthur,{ }ſchon{ }ſeit 8 Tagen bin ich für \label{K_L04286-1v}\edtext{Mittwoch}{\lemma{\textnormal{\emph{Mittwoch}}}\Cendnote{\textnormal{der 9. 12. 1896; was
                  genau geplant war, ist unklar, doch könnte es sich um eine Einladung zur Vorlesung\eventindex{Frankgasse 1@\textbf{Frankgasse 1}!Private Lesung des Prologs zu Mimi, 15.12.1896@Private Lesung des Prologs zu Mimi, 15.12.1896|pwk}
                  des \emph{Prologs}\pwindex{\textcolor{red}{\textsuperscript{XXXX indx1}}!Prolog [zu Mimi]@\strich\emph{Prolog [zu Mimi]}|pwk} zu \emph{Mimi}\pwindex{\textcolor{red}{\textsuperscript{XXXX indx1}}!Mimi. Schattenbilder aus einem Mädchenleben@\strich\emph{Mimi. Schattenbilder aus einem Mädchenleben}|pwk} handeln, die vielleicht
                  wegen Schwarzkopfs\pwindex{Schwarzkopf, Gustav 7.\,11.\,1853 Wien – 13.\,11.\,1939 ebd.@\textsc{Schwarzkopf, Gustav} (7.\,11.\,1853 Wien – 13.\,11.\,1939 ebd.), \emph{Schriftsteller}|pwk} Absage verschoben wurde und 
                  am darauffolgenden Dienstag, dem 15. 12. 1896 stattfand.}}}\label{K_L04286-1}
               bei Nathorff\pwindex{Nathorff, Eugen 12.\,9.\,1847 Frankfurt (Oder) – 2.\,10.\,1902 Wien@\textsc{Nathorff, Eugen} (12.\,9.\,1847 Frankfurt (Oder) – 2.\,10.\,1902 Wien), \emph{Bankier}|pw} zu einer Tarock-Partie
               eingeladen, ich habe auch Ihre Herrn Waſſerma{\geminationn}\pwindex{Wassermann, Jakob 10.\,3.\,1873 Fürth – 1.\,1.\,1934 Altaussee@\textsc{Wassermann, Jakob} (10.\,3.\,1873 Fürth – 1.\,1.\,1934 Altaussee), \emph{Schriftsteller}|pw}, der{ }ſo freundlich war, mir Mitteilung zu
               machen, mein Bedauern ausgeſprochen, daß ich der Vorleſung\eventindex{Frankgasse 1@\textbf{Frankgasse 1}!Private Lesung des Prologs zu Mimi, 15.12.1896@Private Lesung des Prologs zu Mimi, 15.12.1896|pwv} nicht beiwohnen ka{\geminationn}.
               Nun {\pb}danke ich noch Ihnen für Ihre
               freundliche Einladung zum Nachtmahl, die ich ja auch nicht annehmen ka{\geminationn} und grüße Sie herzlichſt.\pend
           
\pstart
           Ihr{\\[\baselineskip]}\spacefill\mbox{Gustav.}\pend
           \leftskip=0em{}\selectlanguage{ngerman}\endnumbering\briefempfaengerindex{Schnitzler, Arthur@\textsc{Schnitzler, Arthur}!zzzSchwarzkopf, Gustav@\emph{von Gustav Schwarzkopf}!1896-12-061@{6. 12. 1896}|)be}\mylabel{L04286h}
\begin{anhang}
\end{anhang}\newcommand{\dateiname}{L04286}\newcommand{\titel}{Gustav Schwarzkopf an Arthur Schnitzler, 6. 12. 1896}\newcommand{\editorInnen}{Herausgegeben von Jahnke, SelmaMüller, Martin Anton}\input{../tex-inputs/latex-leseansicht-abspann}
